\section*{अध्याय \ref{ch:b3:electrode-kinetics} --- इलेक्ट्रोड बलगतिकी और विद्युतविश्लेषण}

\begin{solution}{exo:b3:electrode-kinetics:1}
कैथोडी ढलान $2.303RT/\alpha F = \qty{118}{mV}$ से $\alpha = 59.2/118 = 0.50$ मिलता है; $j_0 = \qty{2.0e-6}{A.cm^{-2}}$,
$\eta = 0$ पर अंतःखंड।
\end{solution}

\begin{solution}{exo:b3:electrode-kinetics:2}
$i_0 = \qty{2.0e-4}{A}$; $R_{\mathrm{ct}} = RT/(Fi_0) = 0.02569/\num{2.0e-4} = \qty{128}{\ohm}$।
\end{solution}

\begin{solution}{exo:b3:electrode-kinetics:3}
\ce{MgSO4} के लिए $I = \frac12(4c + 4c) = 4c$। \qty{1.0}{mmol/L} पर $I = \qty{4.0}{mol.m^{-3}}$ और $\kappa^{-1} =
\qty{0.96}{nm} \times \sqrt{100/4} = \qty{4.8}{nm}$; \qty{0.10}{mol/L} पर $I = \qty{400}{mol.m^{-3}}$, $\kappa^{-1} = \qty{0.48}{nm}$।
\end{solution}

\begin{solution}{exo:b3:electrode-kinetics:4}
$i_{\mathrm c} = C_{\mathrm{dl}}Av = \num{20e-6} \times 0.0707 \times 0.1 = \qty{0.14}{\micro A}$, शिखर के 1\,\% से कम।
\qty{10}{V/s} पर यह \qty{14}{\micro A} है, जबकि शिखर केवल $19\sqrt{100} = \qty{190}{\micro A}$ तक बढ़ता है:
आवेशन धारा $v$ के अनुसार बढ़ती है, फ़ैराडेइक धारा $\sqrt v$ के अनुसार।
\end{solution}

\begin{solution}{exo:b3:electrode-kinetics:5}
$D = \pi\bigl(i\sqrt t/nFAc^*\bigr)^2 = \pi(\num{12.2e-6}/(96\,485 \times 0.0707 \times \num{1.00e-6}))^2 =
\qty{1.0e-5}{cm^2.s^{-1}}$।
\end{solution}

\begin{solution}{exo:b3:electrode-kinetics:6}
$FAc^* = \qty{1.364e-2}{C.cm^{-1}}$, $F/RT = \qty{38.92}{V^{-1}}$। \qty{50}{mV/s} पर: $\sqrt{38.92 \times 0.05 \times \num{7.0e-6}} =
\num{3.69e-3}$, $i_{\mathrm p} = 0.4463 \times \num{1.364e-2} \times \num{3.69e-3} = \qty{22.5}{\micro A}$; \qty{500}{mV/s} पर $\sqrt{10}$ गुना
अधिक, \qty{71}{\micro A}।
\end{solution}

\begin{solution}{exo:b3:electrode-kinetics:7}
(a) उत्क्रमणीय। (b) \omterm{def:b3:electrode-kinetics:reversibility}{अर्ध-उत्क्रमणीय}: पृथक्करण क्रमवीक्षण दर के साथ बढ़ता है।
(c) इलेक्ट्रॉन स्थानांतरण के बाद एक रासायनिक चरण आता है (EC): धीमे क्रमवीक्षणों पर
उत्पाद वापसी क्रमवीक्षण से पहले उपभुक्त हो जाता है; तेज़ क्रमवीक्षण रसायन से आगे निकल जाता है।
\end{solution}

\begin{solution}{exo:b3:electrode-kinetics:8}
$K_{ij}a_j/a_i = \num{2e-4} \times 0.14/\num{4.0e-3} = 0.007$: इलेक्ट्रोड 0.7\,\% अधिक पढ़ता है।
\end{solution}

\begin{solution}{exo:b3:electrode-kinetics:9}
न्यूनतम वर्ग: $b = \qty{0.1093}{\micro A.L.\micro g^{-1}}$, $a = \qty{0.466}{\micro A}$, $s = \qty{0.011}{\micro A}$; $c_0 = a/b =
\qty{4.26}{\micro g.L^{-1}}$, जहाँ $u(c_0) = \frac sb\sqrt{\frac1n + \frac{\bar y^2}{b^2S_{xx}}} = \qty{0.12}{\micro g.L^{-1}}$
($n = 4$, $\bar y = 1.286$, $S_{xx} = 125$)।
\end{solution}

\begin{solution}{exo:b3:electrode-kinetics:10}
$Q = 0.0500 \times 1800 = \qty{90.0}{C}$; $n(\ce{Cu}) = Q/2F = \qty{4.66e-4}{mol}$, \qty{29.6}{mg}।
परिणाम केवल आवेश पर, जो एक धारा और एक समय से मापा जाता है, और
फ़ैराडे स्थिरांक पर निर्भर करता है: किसी मानक की आवश्यकता नहीं, बशर्ते वैद्युत-अपघटन
पूर्ण हो और उसकी धारा दक्षता 100\,\% हो।
\end{solution}

\begin{solution}{exo:b3:electrode-kinetics:11}
$\sqrt{DFv/RT} = \sqrt{\num{1.0e-5} \times 38.92 \times 0.1} = \qty{6.24e-3}{cm.s^{-1}}$: $\Lambda = 1.6$, \omterm{def:b3:electrode-kinetics:reversibility}{अर्ध-उत्क्रमणीय}।
\qty{10}{V/s} पर $\Lambda = 0.16$, अब भी \omterm{def:b3:electrode-kinetics:reversibility}{अर्ध-उत्क्रमणीय} परंतु उत्क्रमणीय सीमा से
और दूर; तरंग केवल लगभग \qty{1}{mV/s} से नीचे ही उत्क्रमणीय होगी।
\end{solution}

\begin{solution}{exo:b3:electrode-kinetics:12}\emergencystretch=3em
$s/b = \qty{0.0838}{mM}$ और $b^2S_{xx} = \qty{203.5}{\micro A^2}$। $\bar y$ पर: $0.0838\sqrt{1 + 1/7} = \qty{0.090}{mM}$।
\qty{18.0}{\micro A} पर: $0.0838\sqrt{1.143 + 83.0/203.5} = \qty{0.104}{mM}$। $\bar y$ पर तीन पाठ्यांक:
$0.0838\sqrt{1/3 + 1/7} = \qty{0.058}{mM}$।
\end{solution}

\begin{solution}{pb:b3:electrode-kinetics:1}
\textbf{1.} $\ce{C6H12O6 -> C6H10O6 + 2 H+ + 2 e-}$ और $\ce{[Fe(CN)6]^3- + e- -> [Fe(CN)6]^4-}$;
समग्र
\[
  \ce{C6H12O6 + 2 [Fe(CN)6]^3- -> C6H10O6 + 2 [Fe(CN)6]^4- + 2 H+} .
\]
इलेक्ट्रोड पर, $\ce{[Fe(CN)6]^4- -> [Fe(CN)6]^3- + e-}$।
\textbf{2.} दो।
\textbf{3.} रक्त में घुली ऑक्सीजन कम और परिवर्ती होती है, और उसका उत्पाद,
हाइड्रोजन पेरॉक्साइड, केवल उच्च विभव पर ऑक्सीकृत होता है जहाँ अनेक अन्य
स्पीशीज़ भी अभिक्रिया करती हैं; मध्यस्थ एक ज्ञात आधिक्य में उपस्थित रहता है।
\textbf{4.} ताकि उसकी पृष्ठीय सांद्रता शून्य हो: तब धारा केवल विसरण द्वारा
सीमित होती है और विभव के छोटे परिवर्तनों के प्रति असंवेदी रहती है।
\textbf{5.} धारा $t^{-1/2}$ के अनुसार गिरती है; पाठ्यांक उसी समय पर लिए जाने चाहिए जो
अंशांकन में प्रयुक्त हुआ।
\textbf{6.} कॉटरेल: $i\sqrt t = 42.0$, 41.4, 41.7, 41.6, 41.8: 1\,\% के भीतर स्थिर।
\textbf{7.} $D = \pi(\text{ढलान}/FAc)^2 = \pi(\num{41.8e-6}/(96\,485 \times 0.0300 \times \num{1.00e-5}))^2 = \qty{6.55e-6}{cm^2.s^{-1}}$।
\textbf{8.} $\sqrt{\pi \times \num{6.55e-6} \times 5} = \qty{0.010}{cm} = \qty{100}{\micro m}$: अवक्षय लगभग \qty{5}{s} पर
परत के शीर्ष तक पहुँच जाता है; बाद के पाठ्यांक कॉटरेल नियम से नीचे गिरेंगे।
\textbf{9.} आधी: \qty{9.35}{\micro A}।
\textbf{10.} पृष्ठभूमि धारा: व्यतिकारियों और अशुद्धियों का ऑक्सीकरण,
आवेशन, मध्यस्थ का अपना अवशिष्ट हेक्सासायनोफ़ेरेट(II)।
\textbf{11.} हाँ: अवशिष्ट, अधिकतम \qty{0.095}{\micro A}, बिना किसी प्रवृत्ति के बिखरे हैं। उच्च ग्लूकोस पर मध्यस्थ, जो
सीमित मात्रा में उपस्थित है, या \omterm{def:b3:complex-kinetics:enzyme}{एंज़ाइम} संतृप्त हो जाता है और रेखा मुड़ जाती है।
\textbf{12.} $(7.10 - 0.356)/0.9189 = \qty{7.34}{mM}$।
\textbf{13.} $0.0838\sqrt{1 + 1/7 + (7.10 - 8.89)^2/203.5} = 0.0838 \times 1.076 = \qty{0.090}{mM}$।
\textbf{14.} $1/m$ पद (एकल पाठ्यांक); बार-बार पाठ्यांक, या कई पट्टियाँ,
इसे घटाती हैं।
\textbf{15.} $3 \times 0.077/0.919 = \qty{0.25}{mM}$।
\textbf{16.} $7.34 \times 180.16 = \qty{1322}{mg.L^{-1}} = \qty{132}{mg.dL^{-1}}$।
\textbf{17.} यह अपनी ऑक्सीकरण धारा जोड़ देता है: मीटर अधिक पढ़ता है।
\textbf{18.} ग्लूकोस-रहित पट्टी का एक \omterm{def:b3:electrode-kinetics:cv}{वोल्टामोग्राम}, जहाँ हेक्सासायनोफ़ेरेट(II) ऑक्सीकृत होता है वहाँ,
ऐस्कॉर्बेट की एक ऐनोडी तरंग दिखाता है।
\textbf{19.} निम्नतर विभव पर कम पदार्थ ऑक्सीकृत होते हैं।
\textbf{20.} अंशांकन रेखा लागू करने से पहले रिक्त धारा को \omterm{def:b3:complex-kinetics:enzyme}{एंज़ाइम} इलेक्ट्रोड के पाठ्यांक से
घटा दिया जाता है।
\textbf{21.} $D$ प्रति केल्विन कुछ प्रतिशत बढ़ता है और धारा $\sqrt D$ के अनुसार: मीटर
ताप मापते हैं और उसके लिए संशोधन करते हैं।
\textbf{22.} रक्त की श्यानता और लाल-कोशिका अंश $D$ और धारा को बदलते हैं;
समान आधात्री में बने मानक इन प्रभावों को निरस्त कर देते हैं।
\textbf{23.} पट्टी-दर-पट्टी विचरण (इलेक्ट्रोड क्षेत्रफल, \omterm{def:b3:complex-kinetics:enzyme}{एंज़ाइम} और मध्यस्थ की मात्रा),
ताप और रक्त संघटन अंशांकन की अपनी अनिश्चितता में जुड़ जाते हैं।
\textbf{24.} \textbf{$c(\text{ग्लूकोस}) = 7.34 \pm \qty{0.09}{mmol.L^{-1}}$} (मानक अनिश्चितता), लगभग \qty{132}{mg.dL^{-1}}।
\end{solution}
